
\documentclass[
	% -- opções da classe memoir --
	12pt,				% tamanho da fonte
	openright,			% capítulos começam em pág ímpar (insere página vazia caso preciso)
	oneside,			% para impressão em verso e anverso. Oposto a oneside
	a4paper,			% tamanho do papel. 
	% -- opções da classe abntex2 --
	%chapter=TITLE,		% títulos de capítulos convertidos em letras maiúsculas
	%section=TITLE,		% títulos de seções convertidos em letras maiúsculas
	%subsection=TITLE,	% títulos de subseções convertidos em letras maiúsculas
	%subsubsection=TITLE,% títulos de subsubseções convertidos em letras maiúsculas
	% -- opções do pacote babel --
	english,			% idioma adicional para hifenização
	brazil				% o último idioma é o principal do documento
	]{abntex2}

% ---
% Pacotes básicos 
% ---
\usepackage{unitins}
\usepackage{lmodern}			% Usa a fonte Latin Modern			
\usepackage[T1]{fontenc}		% Selecao de codigos de fonte.
\usepackage[utf8]{inputenc}		% Codificacao do documento (conversão automática dos acentos)
\usepackage{lastpage}			% Usado pela Ficha catalográfica
\usepackage{indentfirst}		% Indenta o primeiro parágrafo de cada seção.
\usepackage{color}				% Controle das cores

\usepackage{graphicx}			% Inclusão de gráficos
% \usepackage{svg}
% \usepackage{xcolor}
\usepackage{subfig}
\usepackage{microtype} 			% para melhorias de justificação
\usepackage{soul}
\usepackage{amssymb}
\usepackage{amsmath}
\usepackage{spreadtab}
\usepackage{multirow}
\usepackage{amsthm}
\usepackage{url}
\usepackage[portuguese, ruled, linesnumbered]{algorithm2e}
\usepackage{float}
\usepackage{appendix}
\usepackage{fancyhdr}
\usepackage{lastpage}
\usepackage{verbatim}

% ---

% ---
% Pacotes de citações
% ---
\usepackage{hyperref}

%\usepackage[brazilian,hyperpageref]{backref}	 % Paginas com as citações na bibl
\usepackage[alf,abnt-etal-list=2 ]{abntex2cite}	% Citações padrão ABNT


\usepackage[table]{xcolor}
\definecolor{lightgray}{gray}{0.9}
\graphicspath{{imagens/}}

\theoremstyle{theorem}
\newtheorem{teo}{Teorema}[chapter]
\newtheorem{lema}[teo]{Lema}
\theoremstyle{definition}
\newtheorem{defi}[teo]{Definição}

\usepackage{svg}
\svgpath{{imagens/svgimgs}} % <- using \svgpath to avoid warning
\usepackage{amsmath}
\usepackage{url}


% ---
% Informações de dados para CAPA e FOLHA DE ROSTO
% ---

\titulo{Governança de TI em Municípios de Baixa Maturidade: proposta de um guia mínimo para elevação do IGovTI nas gestões municipais do Tocantins}
\autor{John Helder Cardoso Alves}
\local{Palmas}
\data{2026}
\orientador{Prof. Me. Tamirys Virgulino Ribeiro Prado}
\instituicao{%
  Fundação Universidade do Tocantins
  \par
  Curso de Sistemas de Informação
  }
\tipotrabalho{Trabalho de Conclusão de Curso (Graduação)}
% O preambulo deve conter o tipo do trabalho, o objetivo, 
% o nome da instituição e a área de concentração 
\preambulo{
Trabalho apresentado ao Curso de Sistemas de Infor-
mação da Fundação Universidade do Tocantins -
UNITINS como parte dos requisitos para a obtenção
do grau de Bacharel em Sistemas de Informação,
sob a orientação da professora Me. Tamirys Virgulino Ribeiro Prado
}

% ---


% ---
% Configurações de aparência do PDF final

% alterando o aspecto da cor azul
\definecolor{blue}{RGB}{41,5,195}

% informações do PDF
\makeatletter
\hypersetup{
     	%pagebackref=true,
		pdftitle={\@title}, 
		pdfauthor={\@author},
    	pdfsubject={\imprimirpreambulo},
	    pdfcreator={LaTeX with abnTeX2},
		pdfkeywords={Algoritmo}{trabalho acadêmico}, 
		colorlinks=true,       		% false: boxed links; true: colored links
    	linkcolor=blue,          	% color of internal links
    	citecolor=blue,        		% color of links to bibliography
    	filecolor=magenta,      		% color of file links
		urlcolor=blue,
		bookmarksdepth=4
}
\makeatother
% --- 

% --- 
% Espaçamentos entre linhas e parágrafos 
% --- 

% O tamanho do parágrafo é dado por:
\setlength{\parindent}{1.3cm}

% Controle do espaçamento entre um parágrafo e outro:
\setlength{\parskip}{0.2cm}  % tente também \onelineskip

% ---
% compila o indice
% ---
\makeindex
% ---

% ----
% Início do documento
% ----
\begin{document}

% Retira espaço extra obsoleto entre as frases.
\frenchspacing 

% ----------------------------------------------------------
% ELEMENTOS PRÉ-TEXTUAIS
% ----------------------------------------------------------
% \pretextual
% ----------------------------------------------------------

% ----------------------------------------------------------
% ELEMENTOS PRÉ-TEXTUAIS
% ----------------------------------------------------------
% \pretextual

% ---
% Capa
% ---
\imprimircapa
% ---

% ---
% Folha de rosto
% (o * indica que haverá a ficha bibliográfica)
% ---
\imprimirfolhaderosto
% ---

% ---
% Inserir folha de aprovação
% ---

% Isto é um exemplo de Folha de aprovação, elemento obrigatório da NBR
% 14724/2011 (seção 4.2.1.3). Você pode utilizar este modelo até a aprovação
% do trabalho. Após isso, substitua todo o conteúdo deste arquivo por uma
% imagem da página assinada pela banca com o comando abaixo:
%
% \includepdf{folhadeaprovacao_final.pdf}
%
\begin{folhadeaprovacao}

  	
  	
  	\begin{center}
  		\includegraphics[width=1\textwidth]{imagens/unitins.png}
  		\ABNTEXchapterfont\Large   CURSO DE SISTEMAS DE INFORMA{\c{C}}{\~{A}}O
  		
  		\par
  		\vspace*{1cm}     
  		{\ABNTEXchapterfont\bfseries\large \expandafter\MakeUppercase{\imprimirtitulo}  \vspace*{1cm}    }
  		\par
  		{\large \expandafter\MakeUppercase{\imprimirautor}}
  		%\vspace*{\fill}
  		\par
  		\vspace*{1cm}     
  		\hspace{.45\textwidth}
  		\begin{minipage}{.5\textwidth}
  			\small\imprimirpreambulo
  			
  		\end{minipage}%
  	%	\vspace*{\fill}
  	\end{center}
 
  
        
  
   \assinatura{\textbf{\imprimirorientador} \\ Orientador} 
   \assinatura{\textbf{Convidado} \\ Convidado 1}
   \assinatura{\textbf{Convidado} \\ Convidado 2}
   %\assinatura{\textbf{Professor} \\ Convidado 3}
   %\assinatura{\textbf{Professor} \\ Convidado 4}
      
   \begin{center}
    \vspace*{0.5cm}
    {\large\imprimirlocal}
    \par
    {\large\imprimirdata}
  \end{center}

% \begin{center}
%     \centering
%     \includegraphics[scale=0.8]{ficha.pdf} 
% \end{center}

% \begin{center}
%     \centering
%     \includegraphics[scale=0.8]{ata_assinada.pdf} 
% \end{center}

\end{folhadeaprovacao}
% ---

% ---
% Dedicatória
% ---
% \begin{dedicatoria}

% \end{dedicatoria}
% ---

% ---
% Agradecimentos
% --- Es
% \begin{agradecimentos}

% \end{agradecimentos}
% ---

% ---
% Epígrafe
% ---
\begin{epigrafe}
    \vspace*{\fill}
	\begin{flushright}
		\textit{``Não vos amoldeis às estruturas deste mundo, \\
		mas transformai-vos pela renovação da mente, \\
		a fim de distinguir qual é a vontade de Deus: \\
		o que é bom, o que Lhe é agradável, o que é perfeito.”\\
		(Bíblia Sagrada, Romanos 12, 2)}
	\end{flushright}
\end{epigrafe}
% ---

% ---
% RESUMOS
% ---

% resumo em português
\setlength{\absparsep}{18pt} % ajusta o espaçamento dos parágrafos do resumo
\begin{resumo}
Lorem Ipsum is simply dummy text of the printing and typesetting industry. Lorem Ipsum has been the industry's standard dummy text ever since 1966, when designers at Letraset and James Mosley, the librarian at St Bride Printing Library in London, took a 1914 Cicero translation and scrambled it to make dummy text for Letraset's Body Type sheets. It has survived not only many decades, but also the leap into electronic typesetting, remaining essentially unchanged. It was popularised thanks to these sheets and more recently with desktop publishing software like Aldus PageMaker and Microsoft Word including versions of Lorem Ipsum.
Lorem Ipsum is simply dummy text of the printing and typesetting industry. Lorem Ipsum has been the industry's standard dummy text ever since 1966, when designers at Letraset and James Mosley, the librarian at St Bride Printing Library in London, took a 1914 Cicero translation and scrambled it to make dummy text for Letraset's Body Type sheets. It has survived not only many decades, but also the leap into electronic typesetting, remaining essentially unchanged. It was popularised thanks to these sheets and more recently with desktop publishing software like Aldus PageMaker and Microsoft Word including versions of Lorem Ipsum.
Lorem Ipsum is simply dummy text of the printing and typesetting industry. Lorem Ipsum has been the industry's standard dummy text ever since 1966, when designers at Letraset and James Mosley, the librarian at St Bride Printing Library in London, took a 1914 Cicero translation and scrambled it to make dummy text for Letraset's Body Type sheets. It has survived not only many decades, but also the leap into electronic typesetting, remaining essentially unchanged. It was popularised thanks to these sheets and more recently with desktop publishing software like Aldus PageMaker and Microsoft Word including versions of Lorem Ipsum.

\textbf{Palavras-chaves}: Lorem, Ipsum, is e simply.
 
\end{resumo}

% resumo em inglês
\begin{resumo}[Abstract]
 \begin{otherlanguage*}{english}
Lorem Ipsum is simply dummy text of the printing and typesetting industry. Lorem Ipsum has been the industry's standard dummy text ever since 1966, when designers at Letraset and James Mosley, the librarian at St Bride Printing Library in London, took a 1914 Cicero translation and scrambled it to make dummy text for Letraset's Body Type sheets. It has survived not only many decades, but also the leap into electronic typesetting, remaining essentially unchanged. It was popularised thanks to these sheets and more recently with desktop publishing software like Aldus PageMaker and Microsoft Word including versions of Lorem Ipsum.
Lorem Ipsum is simply dummy text of the printing and typesetting industry. Lorem Ipsum has been the industry's standard dummy text ever since 1966, when designers at Letraset and James Mosley, the librarian at St Bride Printing Library in London, took a 1914 Cicero translation and scrambled it to make dummy text for Letraset's Body Type sheets. It has survived not only many decades, but also the leap into electronic typesetting, remaining essentially unchanged. It was popularised thanks to these sheets and more recently with desktop publishing software like Aldus PageMaker and Microsoft Word including versions of Lorem Ipsum.
Lorem Ipsum is simply dummy text of the printing and typesetting industry. Lorem Ipsum has been the industry's standard dummy text ever since 1966, when designers at Letraset and James Mosley, the librarian at St Bride Printing Library in London, took a 1914 Cicero translation and scrambled it to make dummy text for Letraset's Body Type sheets. It has survived not only many decades, but also the leap into electronic typesetting, remaining essentially unchanged. It was popularised thanks to these sheets and more recently with desktop publishing software like Aldus PageMaker and Microsoft Word including versions of Lorem Ipsum.



   \vspace{\onelineskip}
 
   \noindent 
   \textbf{Key-words}: Lorem, Ipsum, is e simply.
 \end{otherlanguage*}
\end{resumo}


% ---
% inserir lista de ilustrações
% ---
\pdfbookmark[0]{\listfigurename}{lof}
\listoffigures*
% ---

% ---
% inserir lista de tabelas
% ---
\newpage
\pdfbookmark[0]{\listtablename}{lot}
\newpage
\listoftables*
\cleardoublepage
% ---

% ---
% inserir lista de abreviaturas e siglas
% ---
\begin{siglas}
  
\end{siglas}
% ---


% ---
% inserir o sumario
% ---
\pdfbookmark[0]{\contentsname}{toc}
\tableofcontents*
\cleardoublepage
% ---



% ELEMENTOS TEXTUAIS
% ---------------------------------[!htb]-------------------------
%\textual

%Capitulos


\mainmatter
\chapter{Introdução}\label{introducao}
\section{Objetivos}
\subsection{Objetivos Específicos}
\section{Justificativa}
\chapter{Referencial Teórico}
O objetivo desta seção é fundamentar teoricamente a pesquisa, delimitando os conceitos, princípios e 
parâmetros institucionais que norteiam o estudo. Este capítulo visa conferir rigor metodológico e 
consistência analítica à avaliação da gestão pública municipal e às suas ferramentas de governança.

\section{Índice de Efetividade da Gestão Municipal}
Concebido originalmente em 2014 pelo Tribunal de Contas do Estado de São Paulo (TCE-SP), o Índice de 
Efetividade da Gestão Municipal (IEGM) foi posteriormente incorporado e aperfeiçoado pelo Instituto 
Rui Barbosa (IRB) - entidade voltada ao desenvolvimento de estudos e pesquisas para o sistema de 
controle externo \cite{iegm_ceara_2022}. Com o intuito de uniformizar a avaliação das administrações 
públicas locais, o índice foi selecionado em 2016 pela Rede Nacional de Indicadores Públicos 
(Indicon) para aplicação em âmbito nacional \cite{iegm_ceara_2022}.

O indicador foi projetado para mensurar o nível de eficácia e desenvolvimento das 
administrações municipais brasileiras. O IEGM engloba sete áreas fundamentais: educação, saúde, 
gestão fiscal, planejamento, meio ambiente, defesa civil e governança de TI \cite{lima_influencia_2023}.

O resultado geral do IEGM varia de 0 a 100, no qual a pontuação máxima representa a aderência 
plena do município às normas e aos controles estabelecidos. Para o cálculo do índice consolidado, 
aplica-se uma métrica de ponderação que reflete a relevância de cada dimensão específica. Assim, a 
composição do IEGM é estruturada pelas seguintes proporções:
\cite{lima_influencia_2023}.

\begin{itemize}
  \item Educação (i-Educ, 20\%)
  \item Saúde (i-Saúde, 20\%)
  \item Gestão Fiscal (i-Fiscal, 20\%)
  \item Meio Ambiente (i-Amb, 10\%)
  \item Proteção dos Cidadãos/Defesa Civil (i-Cidade, 5\%)
  \item Governança de TI (i-Gov TI, 5\%)
\end{itemize}

O Índice conta com a adesão de todos os 27 Tribunais de Contas dos estados e do Distrito Federal, 
consolidando-se como um instrumento nacional de avaliação das administrações públicas locais \cite{araujo_gestoes_municipais_2021}. 
A operacionalização do índice ocorre por meio do Termo de Adesão à Rede de Indicadores Nacionais, 
firmado entre o Instituto Rui Barbosa (IRB) e os órgãos de contas estaduais e distrital. A partir 
desse acordo, as informações prestadas pelos municípios são coletadas e posteriormente validadas 
in loco ou auditadas pelos respectivos tribunais, garantindo a consistência dos dados do índice \cite{araujo_gestoes_municipais_2021}.

\section{Governança de TI na Administração Pública}

\section{Governança Frugal de TI em Municípios de Baixa Maturidade}

\section{Design Science Research(DSR) como Abordagem Metodológica}

\section{Trabalhos Relacionados}




\mainmatter
\chapter{Metodologia da Pesquisa}\label{metodologia}





\mainmatter
\chapter{Resultados e Discussão}\label{resultados}





\mainmatter
\chapter{Considerações Finais}\label{conclusao}






% ----------------------------------------------------------
% ELEMENTOS PÓS-TEXTUAIS
% ----------------------------------------------------------
\postextual
% ----------------------------------------------------------

% ----------------------------------------------------------
% Referências bibliográficas
% ----------------------------------------------------------
\bibliography{bibliografia}

% \include{apendices}
% ----------------------------------------------------------
% Glossário
% ----------------------------------------------------------
%
% Consulte o manual da classe abntex2 para orientações sobre o glossário.
%
%\glossary

%---------------------------------------------------------------------
% INDICE REMISSIVO
%---------------------------------------------------------------------
\phantompart
\printindex
%---------------------------------------------------------------------

\end{document}
\grid
